%% Section 03 — Background and Literature Review
\chapter{Background and Literature Review}
\label{ch:litreview}

\section{Introduction}
\label{sec:lr-intro}
The Antarctic Treaty System (ATS) is a regime negotiated between a set of parties
called consultative parties \textcite{ATS2023RoP}. This group of countries, together with
a second group lacking decision-making power (non-consultative parties), meets
 yearly at the Antarctic Treaty Consultative Meeting (ATCM). Before each meeting,
culsultative parties may submit Working Papers (WP) and Information Papers (IP) up
to 45 days before the ATCM starts \textcite{ATS2023RoP}.
During ATCMs the parties discuss Working Papers and produce outcomes only by consensus.
These outcomes take the form of Measures, Decisions, or Resolutions \textcite{ATCM1995dec1}.
The consensus requirement makes the ATS a notable example of international cooperation,
 but is also its main source of fragility.
As international relations strain and Antarctic environmental pressures intensify,
the risk of losing the capacity to act increases. In 2022, for the first time, the
ATCM parties failed to reach consensus on the final report; even as activity has
increased, the regime has become less responsive to challenges \textcite{Gardiner2025}.


In a sistem based on consensus is or utmost importance to uderstand why it succeds or fails,
And which parties drive it or block it, respectively. Existing literature has described the ATS
by rates and diversity of outputs, and stucturally by describing who co-aouthors with whom 
\textcite{Botelho_CoalitionDynamics,Kumar_ConsistentPatterns}. But there is no study where the
a reconstruction of a directional party intention was reconstructed from the textual records
 on each issue over time.
 Even further, that type of measure has not been used to model the strategic behaviour that 
 produces or breaks consensus. This thesis objective is to measure that and to build that 
 model and asks if they contain the strategic intentions made by the consultative parties.
 It then asks whether these can model historical consensus dynamics.

The approach is a computational pipeline. First, it uses embedding techniques to discover topics, 
and NLI (Natural Language Inference) techniques for stance extraction. These would allow
the construction of the preference measure. Second, it applies NLI to extract the opposition 
or aversion. It then testes the approach by comparing inputs (i.e. working papers) against
 observes outcomes. Finally, a Multi-Agent reinforcement learning is trained based on the
  previous outcomes to simulate the negotiation itself. 

\section{The Antarctic Treaty System as a consensus regime}
\label{sec:ats}
In the ATS, Consultative Parties (CPs) reach outcomes exclusively by consensus. 
These discussions take place at the annual ATCM, where outcomes can take any of the following forms,
 per \textcite[Rule~24]{ATS2023RoP}:
\begin{itemize}
\item \textbf{Measures} are adopted texts containing provisions intended to become legally binding;
 these require approval by the relevant parties' legislative bodies.
\item \textbf{Decisions} are adopted texts concerning internal organisational matters of the ATCM.
\item \textbf{Resolutions} are non-binding recommendations or expressions of intent.
\end{itemize}

There are four types of input:
\begin{itemize}
\item \textbf{Working Papers} (\textcite[Rule~48]{ATS2023RoP}): submitted by CPs and Observers, these require
 discussion at the ATCM and are the main source of outcomes. They are the only input from CPs that can produce
  an outcome; as such, they are the most likely place for each country to state its intentions.
\item \textbf{Secretariat Papers} (\textcite[Rule~49]{ATS2023RoP}): prepared and presented by the Secretariat,
 these produce organisational decisions as output, e.g.\ the budget.
\item \textbf{Information Papers} (\textcite[Rule~50]{ATS2023RoP}): supporting information for Working Papers
 or other discussions. They may be presented by CPs, Observers, Non-Consultative Parties (NCPs), and experts.
\item \textbf{Background Papers} (\textcite[Rule~51]{ATS2023RoP}): submitted by ATCM participants but not 
introduced during the meeting; they exist solely for the record.
\end{itemize}

Of all four, Working Papers are the only input uniquely submitted by CPs, requiring discussion and not confined
 to organisational matters. They are therefore the most likely place for countries to express their strategic 
 intentions within this closed decision-making environment.

The consensus system makes the ATS a notable example of international cooperation; however, the same mechanism 
has become a major challenge as the number of CPs grows. \textcite{Gardiner2025} show that the decision-making 
responsiveness of the ATS is declining; the unprecedented 2022 consensus failure is a clear example.

A notable case of contested consensus occurred in the late 1970s and 1980s, during the negotiations over mineral
resources in Antarctica. The Convention on the Regulation of Antarctic Mineral Resource Activities (CRAMRA) was
produced and then abandoned in favour of the Madrid Protocol (1991), which prohibited all mineral resource 
extraction in Antarctica \parencite{jackson2021antarctica}. This well-documented episode provides a valuable canary
test: a known case with documented positions against which any quantitative stance measure can be validated.


\section{Quantitative and computational study of the ATS}
\label{sec:gametheory}

The study of strategic interaction among self-interested actors was formalised by
\textcite{Nash1951} and later extended to evolutionary and repeated-game settings
\parencite{Axelrod1984, Osborne1994}.
Consensus dynamics in collective decision-making have been modelled through opinion-pooling
and DeGroot-style belief updating \parencite{DeGroot1974}, where agents iteratively revise
their positions toward a weighted average of their neighbours' views.

% TODO: Add coverage of: cooperative game theory (core, Shapley value),
%       veto-player models in political science, coalition formation literature.

\subsection{Affinity and Aversion Dynamics}
\label{sec:affinity-aversion}

% TODO: Develop the theoretical basis for Paper 1 (affinity) and Paper 2 (aversion).
%       Define the formal distinction between positive and negative coalition forces.

\section{Measuring policy preference from political text}
\label{sec:computational-ir}

Agent-based modelling (ABM) offers a bottom-up methodology for studying emergent
collective behaviour from individual rules \parencite{Epstein1996}.
Applications to political and economic systems have grown substantially
\parencite{Farmer2009}, but systematic application to treaty dynamics remains limited.

% TODO: Survey existing computational IR literature, identify the gap this project fills.

\section{Modelling consensus, coalition formation and negotiation}
\label{sec:ai-policy}

Recent work on AI governance \parencite{Dafoe2020} highlights both the transformative
potential and the risks of deploying machine-learning systems in high-stakes institutional
contexts.
This project uses AI methods (reinforcement learning, network analysis) as \emph{tools}
for modelling treaty dynamics rather than as agents within the modelled system.

% TODO: Add specific ML techniques planned — RL for simulating adaptive state strategies,
%       graph neural networks for network topology of treaty relationships.

\section{Summary of Literature Gap}
\label{sec:gap}

% TODO: Synthesise the gap: no formal computational framework exists that:
%       (i) models consensus dynamics in the ATS with empirical calibration,
%       (ii) integrates affinity/aversion forces in a unified agent model, and
%       (iii) connects to broader treaty design implications.
